We need to compute the initial condition for the system for each evaluation of the cost function so that we have a periodic trajectory, 
meaning that the last states of the trajectory we the same as the initial condition.

The solution of the state equation \eqref{cap2-equacao-dinamica} is given by 


\begin{equation}
    x(t) = e^{At} x(0) + \int_{0}^{t}{e^{A(t - \tau)} B u(\tau) d\tau}
\end{equation}

\purple{
    \begin{equation}
        x(t) = e^{At} x(0) + \int_{0}^{t}{e^{A(t - \tau)} \mathbf{b} d\tau}
    \end{equation}
}

We can have different dynamics for each switching instant, so the system behaves as 

\begin{align}
    [0 - t_1]   &: \dot{x}= A_1 x + B_1 u_1 \nonumber \\
    [t_1 - t_2] &: \dot{x}= A_2 x + B_2 u_2 \nonumber \\
    &\qquad \vdots \nonumber \\
    [t_{N-1} - t_N] &: \dot{x}= A_N x + B_N u_N 
\end{align}

We can rewrite the solution for the first region as

\noindent
\begin{align}
    x(t_1) &= e^{A_1 t_1} x(0) + 
            \int_{0}^{t_1}{e^{A_1(t_1 - \tau)} B_1 u_1 d\tau} \nonumber \\
    x(t_1) &= \green{e^{A_1 t_1}} x(0) + 
            \bigg[\orange{\int_{0}^{t_1}{e^{A_1(t_1 - \tau)} B_1 d\tau}}\bigg] u_1 \nonumber \\
    x(t_1) &= \green{F_1} x(0) + \orange{G_1} u_1
\end{align}

\noindent
then the behavior of the other regions of the system is

\noindent
\begin{align}
    x(t_1) &= F_1 x(0) + G_1 u_1 \nonumber \\
    x(t_2) &= F_2 x(1) + G_2 u_2 \nonumber \\
    &\qquad \vdots \nonumber \\ 
    x(t_N) &= F_N x(N-1) + G_N u_N
    \label{cap2-get-x0-multi-equations}
\end{align}

From \eqref{cap2-get-x0-multi-equations} we can establish the relation between $x(t_N)$ and $x(t_1)$

\begin{align}
    \begin{split}
        x(t_N) &= 
        F_N F_{N-1} \ldots F_1 x(0) +
        \green{F_N F_{N-1} \ldots F_2 G_1 x(1) +} \\
        &\quad \green{F_N F_{N-1} \ldots F_2 G_2 x(2) + \ldots +}
        \green{F_N G_{N-1} x(N-1) + G_{N} x(N)}
    \end{split} \nonumber \\
    x(t_N) &= F_N F_{N-1} \ldots F_1 x(0) + \green{c}
\end{align}

We know that $x(t_N) = x(0)$, so the solution for $x(0)$ is

\noindent
\begin{align}
    x(0) = F_N F_{N-1} \ldots F_1 x(0) + c \nonumber \\
    \underbrace{\Big(I - F_N F_{N-1} \ldots F_1\Big)}_{\text{presumed invertible}} x(0) = c \nonumber \\
    x(0) = \Big(I - F_N F_{N-1} \ldots F_1\Big)^{-1} c
\end{align}

\noindent
where
\begin{equation}
    \begin{split}
        c &= F_N F_{N-1} \ldots F_2 G_1 x(1) + \\
        &\quad F_N F_{N-1} \ldots F_2 G_2 x(2) + \ldots +
        F_N G_{N-1} x(N-1) + G_{N} x(N)
    \end{split}
\end{equation}